\documentclass[11pt]{article}
\usepackage[margin=1in]{geometry}
\usepackage{amsmath,amssymb,amsthm}
\usepackage[T1]{fontenc}
\usepackage{lmodern}
\usepackage[hidelinks]{hyperref}
\newtheorem{theorem}{Theorem}
\title{Product stationary laws at capacity two\\A disproof of Conjecture 00000001251}
\author{gaochengzhecpu}
\date{October 9, 2026}
\begin{document}
\maketitle
\begin{abstract}
On the three-site cycle, capacity-two partial exclusion has a full-support
product stationary law with one-site distribution $\operatorname{Bin}(2,1/2)$.
For the alternative convention of unit-rate legal jumps, the product of
three uniform laws on $\{0,1,2\}$ is stationary. Both examples contradict
the assertion that product stationarity requires capacity one.
\end{abstract}

\section{A nondegenerate three-site example}
The source defines $k$-exclusion by the capacity $k$ at each site and asserts
that product stationary states occur only at $k=1$. Consider the three-cycle
with sites $0,1,2$, so every pair of distinct sites is adjacent. Its state
space is $\Omega=\{0,1,2\}^3$. A legal jump $i\to j$ decreases $\eta_i$ by one
and increases $\eta_j$ by one, requiring $\eta_i>0$ and $\eta_j<2$.
Every other coordinate remains unchanged.

First use the partial-exclusion rates
\[
 c_{ij}(\eta)=\eta_i(2-\eta_j),\qquad i\ne j.
\]
The off-diagonal generator $Q(\eta,\xi)$ sums these rates over the legal
jumps producing $\xi$; its diagonal is minus the total exit rate.
Define
\[
 u(0)=\frac14,\quad u(1)=\frac12,\quad u(2)=\frac14,
 \qquad \pi(\eta)=u(\eta_0)u(\eta_1)u(\eta_2).
\]
All 27 states have positive mass, and their masses sum to one.

\begin{theorem}
The law $\pi$ is a product stationary distribution for this capacity-two
process. Moreover,
\[
 \mathbb E_\pi\eta_i=1,\qquad
 \mathbb E_\pi(\eta_0\eta_1\eta_2)=1,
 \qquad
 \mathbb E_\pi(\eta_0\eta_1\eta_2)-\prod_{i=0}^2\mathbb E_\pi\eta_i=0.
\]
\end{theorem}
\begin{proof}
For a legal jump, write $a=\eta_i\in\{1,2\}$ and
$b=\eta_j\in\{0,1\}$. The destination has coordinates $a-1,b+1$.
The binomial weights satisfy
\[
 u(a)u(b)a(2-b)=u(a-1)u(b+1)(b+1)(3-a).
\]
Multiplying by the unaffected site's weight gives detailed balance between
this jump and its reverse. Summing over jumps yields
$\pi(\eta)Q(\eta,\xi)=\pi(\xi)Q(\xi,\eta)$. Thus
\[
 \sum_{\eta\in\Omega}\pi(\eta)Q(\eta,\xi)
 =\pi(\xi)\sum_{\eta\in\Omega}Q(\xi,\eta)=0.
\]
For a finite continuous-time chain, this stationarity equation implies
$\pi e^{tQ}=\pi$ for every $t\geq0$, by the exponential power series.
The three coordinates are independent with mean one, giving all the
stated correlation identities.
\end{proof}

\newpage
\section{The unit-rate convention also has a counterexample}
Another common $k$-exclusion convention assigns rate one to each legal
directed jump. Keep exactly the same state space and legal moves, and call
its generator $Q_U$. Each legal jump has a legal reverse of the same rate,
so $Q_U(\eta,\xi)=Q_U(\xi,\eta)$. Consequently the full-support law
\[
 \pi_U(\eta)=\frac1{27}=\prod_{i=0}^2\frac13
\]
satisfies detailed balance and $\pi_U Q_U=0$. It is the product of three
uniform laws on $\{0,1,2\}$, with capacity $2\ne1$.

Both examples are genuine particle-moving processes. For example,
$(1,0,0)\to(0,1,0)$ has rate two in the first process and rate one in the
second. Neither law is a point mass. Total particle number is conserved;
these are stationary grand-canonical product laws, mixing the invariant
particle-number sectors. The source does not require a fixed total particle
number or an extremal stationary law. Full independence also removes every
mixed three-site correlation deficit that measures failure of factorization.

\section{Formal verification and scope}
Lean uses the finite state space \texttt{Fin 3 $\times$ Fin 3 $\times$ Fin 3},
with actual coordinate occupations, legal single-particle moves, both rate
functions, and both generators. Exact finite computations verify detailed
balance for all pairs of states and both stationarity equations.

\begingroup\raggedright
\begin{itemize}
\setlength{\itemsep}{3pt}
\item \texttt{capacity\_\allowbreak two}, \texttt{weighted\_\allowbreak balance}, and
\texttt{generator\_\allowbreak row\_\allowbreak sum} verify admissibility and balance.
\item \texttt{occupation\_\allowbreak marginals} and \texttt{product\_\allowbreak formula}
identify the normalized binomial product law. The three-point calculation
is \texttt{three\_\allowbreak point\_\allowbreak deficit\_\allowbreak zero}.
\item \texttt{probabilityMass} and \texttt{uniformProbabilityMass} are actual
Mathlib probability mass functions, with full support.
\item \texttt{probabilityMass\_\allowbreak stationary} and
\texttt{uniformProbabilityMass\_\allowbreak stationary} prove their stationarity
equations against the corresponding generators after conversion to real
probabilities.
\end{itemize}
\endgroup
The finite-chain stationarity criterion used is exactly $\pi Q=0$.
Lean verifies that criterion and the generator conditions; it does not
construct a separate space of continuous-time sample paths.
All checked results use only Lean's standard logical axioms. No
unproved placeholders, custom axioms, native evaluation, or external
numerical certificates are used.

\paragraph{Reproduction.}
In \texttt{lean/}, run:
\begin{verbatim}
lake build
lake env lean -DwarningAsError=true Main.lean
\end{verbatim}
Lean 4.19.0 and an exact Mathlib revision are pinned. Rebuild this PDF with
\texttt{tectonic main.tex}. The exact bilingual source and its provenance
are supplied as separate files.

\paragraph{Sources.}
\href{https://github.com/The-Last-Math-Competition/The-Last-Math-Competition/blob/9b795e7a94a6076e49a65a6489e1caf9153abd23/conjectures/00000001251.md}{Conjecture 00000001251};
\href{https://github.com/leanprover-community/mathlib4/commit/c44e0c8ee63ca166450922a373c7409c5d26b00b}{Mathlib c44e0c8ee63c}.
\end{document}
